\documentclass[11pt]{article}
\usepackage[margin=1in]{geometry}
\usepackage{amsmath,amssymb,booktabs,hyperref}
\title{Replication Report: Fermi-Hubbard VQE with the Hamiltonian Variational Ansatz}
\author{Ollie (CherryRd) --- QC-100 replication queue}
\date{2026-06-26 (writeup 2026-07-06)}
\begin{document}
\maketitle

\section*{Target paper}
Cade, Mineh, Montanaro, Stanisic. ``Strategies for solving the Fermi--Hubbard
model on near-term quantum computers.'' \textit{Phys.\ Rev.\ B} \textbf{102},
235122 (2020). arXiv:1912.06007.

\section{Paper summary}
The paper attacks the 2D Fermi--Hubbard ground-state problem on near-term
quantum hardware using VQE with the \emph{Hamiltonian Variational} (HV) ansatz:
alternating layers of $\exp(-i\theta H_{\text{onsite}})$,
$\exp(-i\theta H_{\text{hop},H})$, $\exp(-i\theta H_{\text{hop},V})$ with one
angle per term-group per layer. The reference (depth-0) state is the
non-interacting ($U=0$) ground state projected into the correct particle-number
sector. Central quantitative claims: (i) energy error decreases monotonically
with ansatz depth; (ii) modest depth reaches high accuracy on small lattices;
(iii) parameter count grows linearly with depth (3 params/layer); (iv) HV
exploits model structure far better than generic hardware-efficient ans\"atze.

\section{Scope of this replication}
Classical statevector simulation (\texttt{numpy}, \texttt{scipy}) --- no shots,
no device noise, no hardware. Lattices up to 12 qubits. Ansatz, reference state,
optimizer, and ground-truth exact diagonalization all reimplemented clean-room
without \texttt{openfermion}.

\section{Methods}
\textbf{Hamiltonian:} 1D/2D Fermi--Hubbard with $t=1.0$, $U=2.0$, Jordan--Wigner
mapped, hand-built sparse Pauli operators. \textbf{Lattices:}
$1{\times}2$ (4q), $2{\times}2$ (8q), $1{\times}4$ (8q), $1{\times}6$ (12q),
$2{\times}3$ (12q). \textbf{Ansatz:} HV, per-layer $(\theta_O,\theta_H,\theta_V)$
with term-group evolutions built by eigendecomposition (exact
$e^{-i\theta H_g}$ per group). \textbf{Reference:} $U=0$ ground state in the
target particle-number sector. \textbf{Optimizer:} L-BFGS-B, 3 restarts per
depth. \textbf{Ground truth:} exact lowest eigenvalue in
$(n_\uparrow,n_\downarrow)$ sector.

\section{Results}
Energy error $\Delta E = |E_{\text{VQE}} - E_{\text{exact}}|$ and overlap
fidelity vs.\ exact ground state:
\begin{center}
\begin{tabular}{lrrr}
\toprule
Lattice / depth & $n_{\text{params}}$ & $\Delta E$ & fidelity \\
\midrule
$1{\times}2$, depth 1 & 3  & $\sim2\!\times\!10^{-15}$ & $1.0000$ \\
$2{\times}2$, depth 1 & 3  & $3.3\!\times\!10^{-2}$   & $0.9934$ \\
$2{\times}2$, depth 2 & 6  & $7.7\!\times\!10^{-13}$  & $1.0000$ \\
$1{\times}4$, depth 1 & 3  & $5.0\!\times\!10^{-2}$   & $0.983$  \\
$1{\times}4$, depth 3 & 9  & $9.1\!\times\!10^{-4}$   & $0.9998$ \\
$1{\times}4$, depth 6 & 18 & $1.9\!\times\!10^{-5}$   & $1.0000$ \\
$1{\times}6$, depth 4 & 12 & $8.9\!\times\!10^{-3}$   & $0.995$  \\
$1{\times}6$, depth 8 & 24 & $3.7\!\times\!10^{-4}$   & $0.9998$ \\
$2{\times}3$, depth 4 & 12 & $6.1\!\times\!10^{-3}$   & $0.998$  \\
$2{\times}3$, depth 8 & 24 & $1.6\!\times\!10^{-4}$   & $1.0000$ \\
\bottomrule
\end{tabular}
\end{center}
Depth-monotone convergence to $\sim 10^{-4}$ (chemical-accuracy scale) by
depth 8 on the 12-qubit lattices; machine precision at depth 1--2 on
$1{\times}2$ / $2{\times}2$. Parameter count exactly linear in depth
(3/layer), matching the paper.

\section{Critique (honest)}
\begin{itemize}
  \item \textbf{Headline exercised:} yes for the \emph{strategy claim}. The HV
        ansatz, the reference-state prescription, and the depth-vs-accuracy
        scaling were independently reimplemented and quantitatively reproduced
        against exact diagonalization on five lattices, including two 12-qubit
        cases. The paper's depth-scaling curve character is reproduced.
  \item \textbf{Baseline comparison:} exact diagonalization is used as
        ground-truth per lattice. A cross-check against DMRG or a
        hardware-efficient-ansatz baseline was \emph{not} performed; the
        HV-vs-HEA advantage claim is therefore \emph{unverified} here (only the
        HV-vs-exact side).
  \item \textbf{Ansatz-vs-cost trade-off:} verified quantitatively in the
        depth-vs-error tables above --- linear parameter growth, monotone error
        decrease. This is a direct restatement of the paper's Fig.\ 2-class
        result on our lattice set.
  \item \textbf{Not exercised:} (a) the paper's large-lattice scaling study
        (out of classical statevector reach here); (b) shot-noise / device-noise
        robustness; (c) hardware runs and their measurement-error mitigation ---
        the associated raw device data are not deposited with the paper; (d)
        comparison to alternative structured ans\"atze (number-preserving UCC,
        k-UpCCGSD, ADAPT-VQE).
  \item \textbf{Idealization:} noiseless statevector, exact expectation values,
        no shot budget. Real-device behavior is expected to degrade meaningfully
        at the depths ($\sim 8$) needed for the 12-qubit lattices.
\end{itemize}

\section{Verdict}
\textbf{REPLICATED (strategy).} Coverage 7/10, Agreement 10/10. The core
algorithmic claim (HV ansatz + non-interacting reference state $\Rightarrow$
monotone depth-vs-accuracy convergence to chemical-accuracy on small
Fermi--Hubbard lattices) is fully reproduced. Large-lattice scaling and
hardware/noise sections are out of classical-simulation scope and left as
unexercised.

\end{document}
